\section{Integration}

\subsection{Indefinite Integrals and Techniques}

\begin{definition}
    Functions $F: I \to \mathbb{R}$ such that $F'(x) = f(x)$ for all $x \in (a, b)$ are called \textbf{antiderivatives} of $f$.
\end{definition}

\begin{definition}
   The family of all anti-derivatives of a function $f: I \to \mathbb{R}$ for $I \subseteq \mathbb{R}$ an open interval is called the \textbf{indefinite integral} and is denoted by
   $$\int f(x) dx = \{ F: I \to \mathbb{R} : F'(x) = f(x) \} = F(x) + C $$
\end{definition}

\begin{theorem}
    Let $f, g: I \to \mathbb{R}$ be integrable functions and let $c \in \mathbb{R}$. Then
    $$\int (f(x) \pm g(x)) dx = \int f(x) dx \pm \int g(x) dx$$
    $$\int c f(x) dx = c \int f(x) dx$$
\end{theorem}

\begin{remark}
    Some important basic antiderivatives:
    \begin{itemize}
        \item $f(x) = 0 \quad \implies \quad F(x) = C$
        \item $f(x) = 1 \quad \implies \quad F(x) = x + C$
        \item $f(x) = x^p \quad \implies \quad F(x) = \frac{1}{p+1} x^{p+1} + C \quad (p \neq -1)$
        \item $f(x) = e^x \quad \implies \quad F(x) = e^x + C$
        \item $f(x) = a^x \quad \implies \quad F(x) = \frac{a^x}{\ln(a)} + C$
        \item $f(x) = \sin(x) \quad \implies \quad F(x) = -\cos(x) + C$
        \item $f(x) = \cos(x) \quad \implies \quad F(x) = \sin(x) + C$
        \item $f(x) = \frac{1}{x} \quad \implies \quad F(x) = \ln|x| + C$
        \item $f(x) = \frac{1}{1+x^2} \quad \implies \quad F(x) = \arctan(x) + C$
    \end{itemize}
\end{remark}

\begin{theorem}
    \textbf{Integration by Parts:} For differentiable functions $f(x)$ and $g(x)$ we have
    $$\int f'(x)g(x)dx = f(x)g(x) - \int f(x)g'(x)dx $$
\end{theorem}

\begin{remark}
    The equation is derived from the product rule for differentiation. 
\end{remark}

\begin{example}
    Let $f'(x) = e^x \implies f(x) = e^x$ and $g(x) = x \implies g'(x) = 1$. Then 
    $$\int x e^x dx = x e^x - \int e^x \cdot 1 dx = x e^x - e^x + C$$
\end{example}

\begin{theorem}
    \textbf{Substitution Rule:} For functions $f(x)$ and $g(x)$ we have
    $$\int f(g(x))g'(x)dx  = \int f(u) u' dx = \int f(u) \frac{du}{dx} dx = \int f(u) du = F(g(x)) + C $$
    where $F$ is an antiderivative of $f$.
\end{theorem}

\begin{remark}
    The equation is derived from the chain rule for differentiation. 
\end{remark}

\begin{example}
    Evaluate $\int 2x e^{x^2} dx$. 
    Let $u = x^2 \implies \frac{du}{dx} = 2x \implies du = 2x dx$.
    Substituting gives:
    $$\int e^u du = e^u + C = e^{x^2} + C$$
\end{example}

\subsection{The Definite Integral (Riemann Integral)}

\begin{remark}
    \textbf{Upper and Lower Sums (Stair Functions):} One can partition the integration interval into subintervals. For each subinterval we can approximate the area of the function with the area of a rectangle. The sum of the areas of these rectangles is an approximation of the area under the function. Taking the infimum of the function on each subinterval yields the lower sum $L$. Taking the supremum yields the upper sum $U$.
\end{remark}

\begin{definition}
    Let $U$ and $L$ denote the upper and lower sums of a bounded function $f$. Then $f$ is called \textbf{Riemann-integrable} iff $\inf U = \sup L$. The value of the definite integral $\int_a^b f(x)dx$ is then this common value.
\end{definition}

\begin{theorem}
    \textbf{Integrability Conditions:}
    \begin{itemize}
        \item All continuous functions on $[a,b]$ are Riemann-integrable.
        \item All monotone functions on $[a,b]$ are Riemann-integrable.
    \end{itemize}   
\end{theorem}

\begin{theorem}
    Let $f$ and $g$ be integrable and $c \in \mathbb{R}$ then:
    \begin{itemize}
        \item Linearity: $\int_a^b (f(x) \pm g(x))dx = \int_a^b f(x)dx \pm \int_a^b g(x)dx$
        \item Monotonicity: $f(x) \leq g(x) \implies \int_a^b f(x)dx \leq \int_a^b g(x)dx$
        \item Triangle inequality: $\left| \int_a^b f(x)dx \right| \leq \int_a^b |f(x)|dx$
        \item Inversion of bounds: $\int_{a}^{b} f(x)dx = -\int_{b}^{a} f(x)dx$
        \item Partition of interval: $\int_{a}^{c} f(x)dx = \int_{a}^{b} f(x)dx + \int_{b}^{c} f(x)dx$
    \end{itemize}    
\end{theorem}

\begin{theorem}
    \textbf{Mean Value Theorem for Integrals:}
    Let $f: [a,b] \to \mathbb{R}$ be continuous. Then there exists a $c \in [a,b]$ such that
    $$\int_{a}^{b} f(x)dx = (b-a)f(c)$$
\end{theorem}

\begin{theorem}
    \textbf{Fundamental Theorem of Integral Calculus:}
    Let $f: [a,b] \to \mathbb{R}$ be continuous. 
    Then for all $C \in \mathbb{R}$, $F(x)$ is an antiderivative of $f$ where:
    $$F(x) = \int_{a}^{x} f(y)dy + C$$
\end{theorem}

\begin{proof}
    As $f$ is continuous it is integrable on $[a,x]$ for any $x \in [a,b]$.
    Let $x_0 \in [a,b]$ and $\varepsilon > 0$ arbitrary. Then there exist a $\delta > 0$ such that for any $x \in [a,b]$ with $0 < |x-x_0| < \delta$ we have $|f(x) - f(x_0)| < \varepsilon$. 
    Then for $x \in [a,b]$ with $0 < |x-x_0| < \delta$:
    $$\left| \frac{F(x)-F(x_0)}{x-x_0} - f(x_0) \right| = \left| \frac{1}{x-x_0} \left( \int_{a}^{x} f(y)dy -\int_{a}^{x_0} f(y)dy \right) - f(x_0) \right|$$
    $$= \left| \frac{1}{x-x_0} \int_{x_0}^{x} f(y)dy - \frac{1}{x-x_0} \int_{x_0}^{x} f(x_0) dy \right| = \left| \frac{1}{x-x_0} \int_{x_0}^{x} (f(y)-f(x_0)) dy \right|$$
    $$\leq \frac{1}{|x-x_0|} \left| \int_{x_0}^{x} |f(y)-f(x_0)| dy \right| \leq \frac{1}{|x-x_0|} \left| \int_{x_0}^{x} \varepsilon dy \right| = \varepsilon$$
\end{proof}

\begin{proof}
    \textbf{Uniqueness of Antiderivatives:} Let $G(x)$ be another antiderivative of $f: [a,b] \to \mathbb{R}$ where $F(x) = \int_{a}^{x} f(y)dy + C$. We define $H(x) = F(x) - G(x)$ and get
    $$H'(x) = F'(x) - G'(x) = f(x) - f(x) = 0$$ 
    Hence $H(x)$ must be a constant $C'$ and $F(x) = G(x) + C'$. Thus, antiderivatives are unique up to a constant.
\end{proof}

\subsection{Sequences of Functions and Integrals}

\begin{definition}
    Let $(f_n)_{n \in \mathbb{N}_0}$ be a sequence of functions. 
    \begin{itemize}
        \item \textbf{Pointwise Convergence}: It converges pointwise to $f$ if for every $x$, the sequence of real numbers $(f_n(x))$ converges to $f(x)$.
        \item \textbf{Uniform Convergence}: It converges uniformly to $f$ if for every $\epsilon > 0$, there exists an $N$ such that for all $n \ge N$ and \textbf{for all} $x$ simultaneously, $|f_n(x) - f(x)| < \epsilon$.
    \end{itemize}
\end{definition}

\begin{theorem}
    \textbf{Swapping Limit and Integral:} If a sequence of integrable functions $(f_n)$ converges \textbf{uniformly} to $f$ on $[a,b]$, then the limit function $f$ is also integrable, and we can swap the limit and the integral:
    $$\int_a^b f(x)dx = \lim_{n \to \infty} \int_a^b f_n(x) dx$$
\end{theorem}

\subsection{Definite Integration Techniques}

\begin{lemma}
    If $F: [a,b] \to \mathbb{R}$ is continuously differentiable, then:
    $$F(x) = F(a) + \int_{a}^{x} F'(t)dt \quad \text{and} \quad \int_{a}^{b} f(t)dt = F(b) - F(a) = \Big[ F(x) \Big]_a^b$$
\end{lemma}

\begin{lemma}
    \textbf{Substitution for Definite Integrals:} Let $f$ be continuous and $g$ continuously differentiable.
    $$\int_{a}^{b} f(g(x))g'(x)dx  = \int_{g(a)}^{g(b)} f(u) du$$
    
    \textbf{Version 2:} If $f'(x) \neq 0$ and $f^{-1}$ is the inverse of $f$:
    $$\int_a^b g(f(x)) dx = \int_{f(a)}^{f(b)} g(y) (f^{-1})'(y) dy$$
\end{lemma}

\begin{example}
    Evaluate $\int_0^{\pi/2} \sin(x) \cos(x) dx$.
    Let $u = \sin(x)$. Then $\frac{du}{dx} = \cos(x) \implies du = \cos(x) dx$.
    For definite integrals, we must also substitute the bounds:
    Lower bound: $x = 0 \implies u = \sin(0) = 0$.
    Upper bound: $x = \pi/2 \implies u = \sin(\pi/2) = 1$.
    Substituting gives:
    $$\int_0^1 u \, du = \Big[ \frac{1}{2} u^2 \Big]_0^1 = \frac{1}{2} (1)^2 - \frac{1}{2} (0)^2 = \frac{1}{2}$$
\end{example}

\begin{lemma}
    \textbf{Integration by Parts for Definite Integrals:} Let $f, g$ be continuously differentiable.
    $$\int_{a}^{b} f'(x)g(x)dx  = \Big[ f(x)g(x) \Big]_{a}^{b} - \int_{a}^{b} f(x)g'(x)dx$$
\end{lemma}

\begin{example}
    Evaluate $\int_1^e \ln(x) dx$.
    We can write this as $\int_1^e 1 \cdot \ln(x) dx$.
    Let $f'(x) = 1 \implies f(x) = x$.
    Let $g(x) = \ln(x) \implies g'(x) = \frac{1}{x}$.
    Using integration by parts for definite integrals:
    $$\int_1^e 1 \cdot \ln(x) dx = \Big[ x \ln(x) \Big]_1^e - \int_1^e x \cdot \frac{1}{x} dx$$
    $$= (e \ln(e) - 1 \ln(1)) - \int_1^e 1 dx = e - \Big[ x \Big]_1^e = e - (e - 1) = 1$$
\end{example}
